% !PRINTED 216
\weberpage{216}
%<<CONT
wird, darin kann $\gamma, \delta$ durch $\gamma + h\alpha$, $\delta + h\beta$ ersetzt werden,
und indem man $h$ ein volles Restsystem modulo $n$ durchlaufen
läßt, erkennt man, daß zu jedem Zahlenpaar $\alpha, \beta$ $n$ der Bedingung
(23) genügende Zahlenpaare $\gamma, \delta$ gehören. Demnach ist die Ordnung der
% [sic] The source prints (23); the determinant condition to which the
% pairs gamma, delta must conform is (22).
Gruppe $\mathfrak{B}$:
\[
\nu = n^3 \Pi\left(1 - \frac{1}{p^2}\right),
\]
oder, wenn wir noch die numerische Funktion

\weq{25}{\psi(n) = n\Pi\left(1 + \frac{1}{p}\right)}
einführen:

\weq{26}{\nu = n\varphi(n)\psi(n)}
und die Ordnung der Gruppe $\mathfrak{A}$:

\weq{27}{\mu = n\varphi(n)^2\psi(n).}

\begin{center}
\S\,64. \gesperrt{Die irreduzibeln Faktoren der Teilungsgleichung.}
\end{center}
Ist
\[
\begin{aligned}
\nu &= \partial\mu - b\mu'\\
\nu' &= -c\mu + a\mu'
\end{aligned}
\]
und $a\partial - bc$ relativ prim zu $n$, so ist der größte gemeinschaftliche
Teiler von $\mu, \mu', n$ zugleich der größte gemeinschaftliche
Teiler von $\nu, \nu', n$. Die Wurzeln $x_{\mu,\mu'}$ der Teilungsgleichung
zerfallen also nach dem größten gemeinschaftlichen Teiler von
$\mu, \mu', n$ in Systeme, die durch die Substitutionen der Gruppe $\mathfrak{A}$
immer nur ineinander übergehen, d.\,h. die Gruppe $\mathfrak{G}$ ist \gesperrt{intransitiv}, % Source verified: the group letter is Fraktur G.
und die Systeme der \gesperrt{Intransitivität} sind die
Systeme $x_{\mu,\mu'}$, in denen $\mu, \mu', n$ einen und denselben größten
gemeinschaftlichen Teiler haben. Die Teilungsgleichung ist also
reduzibel (außer wenn $n$ eine Primzahl ist) und die Wurzeln
eines dieser Systeme der Imprimitivität % [sic] The source changes from systems of intransitivity to imprimitivity here.
genügen einer rationalen
Gleichung (Bd.\,I, \S\,157).

Nach dem vorhergehenden Paragraphen gibt es $\varphi(n)\psi(n)$
Zahlenpaare $\mu, \mu'$, deren größter gemeinschaftlicher Teiler mit $n$
gleich 1 ist, und die diesen Zahlenpaaren entsprechenden Wurzeln
$x_{\mu,\mu'}$ genügen daher einer rationalen Gleichung des Grades $\varphi(n)\psi(n)$.
Diese Gleichung nennen wir die \gesperrt{eigentliche Teilungsgleichung}
für den Divisor $n$, weil nur dann, wenn $\mu, \mu', n$ ohne gemeinsamen
Teiler sind, $x_{\mu,\mu'}$ nicht zugleich Wurzel einer Teilungs-
%CONT>>
